% !TeX program = lualatex
\documentclass[a4paper,11pt]{article}
\usepackage{fontspec}
\usepackage{newunicodechar}
\newunicodechar{♭}{\ensuremath{\flat}}
\IfFontExistsTF{TeX Gyre Pagella}{\setmainfont{TeX Gyre Pagella}}{\setmainfont{P052}}
\IfFontExistsTF{TeX Gyre Heros}{\setsansfont{TeX Gyre Heros}}{\setsansfont{Nimbus Sans}}
\IfFontExistsTF{TeX Gyre Cursor}{\setmonofont{TeX Gyre Cursor}[Scale=0.88]}{\setmonofont{Nimbus Mono PS}[Scale=0.88]}
\usepackage{babel}
\babelprovide[import,main]{spanish}
\usepackage{microtype}
\usepackage[a4paper,margin=2.6cm,headheight=23pt,headsep=10pt,footskip=1.1cm]{geometry}
\raggedbottom
\setlength{\emergencystretch}{3em}
\hyphenpenalty=10000
\exhyphenpenalty=10000
\clubpenalty=10000
\widowpenalty=10000
\usepackage{xcolor}
\definecolor{navybg}{HTML}{0B1D2E}
\definecolor{navydeep}{HTML}{123049}
\definecolor{lightblue}{HTML}{8FD3FE}
\definecolor{amber}{HTML}{F5A623}
\definecolor{softgray}{HTML}{5B6570}
\usepackage{titlesec,needspace}
\titleformat{\section}{\sffamily\Large\bfseries\color{navydeep}\raggedright}{\thesection}{10pt}{}
\titleformat{\subsection}{\sffamily\large\bfseries\color{navydeep}\raggedright}{\thesubsection}{8pt}{}
\titlespacing*{\section}{0pt}{20pt}{9pt}
\usepackage{fancyhdr}
\pagestyle{fancy}
\fancyhf{}
\fancyhead[C]{\small\sffamily\color{softgray}Drop Route --- Diseño Técnico del Software}
\fancyfoot[L]{\small\sffamily\color{softgray}Revisión parcial 1.0.2 · 16/10/2026}
\fancyfoot[R]{\small\sffamily\color{softgray}\thepage}
\renewcommand{\headrulewidth}{0.4pt}
\usepackage{array,longtable,booktabs,calc}
\newcounter{none} % Pandoc 3.12 usa \LTcaptype{none} en tablas sin título
\renewcommand{\arraystretch}{1.15}
\usepackage{enumitem}
\setlist[itemize]{leftmargin=1.4em,itemsep=3pt,topsep=5pt}
\usepackage{fancyvrb,framed,seqsplit}
\let\originaltexttt\texttt
\renewcommand{\texttt}[1]{\originaltexttt{\seqsplit{#1}}}
$if(highlighting-macros)$
$highlighting-macros$
$endif$
\providecommand{\tightlist}{\setlength{\itemsep}{2pt}\setlength{\parskip}{0pt}}
\usepackage{graphicx,tikz,xurl}
\usepackage[colorlinks=true,linkcolor=navydeep,urlcolor=navydeep,bookmarksnumbered=true,pdftitle={DTS - Drop Route},pdfauthor={Jesús Quevedo Rodríguez}]{hyperref}
\usepackage{bookmark}
\setlength{\parindent}{0pt}
\setlength{\parskip}{6pt}
\begin{document}
\begin{titlepage}
\newgeometry{margin=0cm}
\thispagestyle{empty}
\begin{tikzpicture}[remember picture,overlay]
\fill[navybg] (current page.south west) rectangle (current page.north east);
\node[anchor=north,text width=16cm,align=center] at ([yshift=-3.3cm]current page.north) {\sffamily\bfseries\fontsize{16}{19}\selectfont\color{amber}DROP ROUTE};
\node[anchor=north,text width=16cm,align=center] at ([yshift=-5cm]current page.north) {\sffamily\bfseries\fontsize{30}{38}\selectfont\color{white}Documento de Diseño\\Técnico del Software};
\node[anchor=north,text width=14cm,align=center] at ([yshift=-9.2cm]current page.north) {\sffamily\fontsize{14}{20}\selectfont\color{lightblue}Voicings de guitarra y rutas II--V--I\\Base de desarrollo para Claude Code};
\begin{scope}[shift={([xshift=2.4cm,yshift=7.7cm]current page.south west)}]
\foreach \y in {0,0.5,1,1.5,2,2.5} {\draw[white,opacity=0.28,line width=0.7pt] (0,\y)--(16,\y);}
\foreach \x in {0,2.5,5,7.5,10,12.5,15} {\draw[lightblue,opacity=0.25] (\x,0)--(\x,2.5);}
\fill[amber] (3.75,1) circle (0.2);
\fill[lightblue] (8.75,1.5) circle (0.2);
\fill[white] (8.75,2) circle (0.2);
\fill[lightblue] (8.75,2.5) circle (0.2);
\end{scope}
\node[anchor=south west,text width=16cm,align=left] at ([xshift=2.4cm,yshift=2cm]current page.south west) {\sffamily\fontsize{11}{17}\selectfont\color{white}\textcolor{amber}{Autor del proyecto:} Jesús Quevedo Rodríguez\\\textcolor{amber}{Revisión parcial:} versión 1.0.2 y lectura por tónica · 16/10/2026\\\textcolor{amber}{Repositorio:} a01653/Drop\_Route\\\textcolor{amber}{Estado:} aplicación de estudio y práctica};
\end{tikzpicture}
\end{titlepage}
\restoregeometry
\tableofcontents
\newpage
$body$
\end{document}
